\documentclass[10pt,a4paper]{amsart}
\usepackage[margin=19mm]{geometry}
\usepackage{amsmath,amssymb,mathtools}
\usepackage[scr=rsfs,cal=euler]{mathalfa}
\usepackage{xfrac}
\usepackage[unicode,hidelinks,bookmarks=false]{hyperref}
\newcommand{\ie}{i.e.\ }
\newcommand{\cf}{cf.\ }
\newcommand{\resp}{respectively\ }
\newcommand{\Subsection}{\subsection}
\providecommand{\nobreakdash}{\nobreak\hbox{-}}
\setlength{\emergencystretch}{2em}
\multlinegap0pt
\usepackage{amssymb}
\usepackage{tikz}
\usepackage{xcolor}
\usepackage{mathtools}
\newtheorem{lemma}{Lemma}
\newcommand{\End}{\operatorname{End}}
\newcommand{\R}{\mathbb R}
\newcommand{\dd}{\,\mathrm d}
\title{A sub-Riemannian endpoint map with no regular values}
\author{Antonio Lerario, Luca Rizzi, Daniele Tiberio}
\date{Source: arXiv:2609.05991v1 (CC BY 4.0)}
\begin{document}
\maketitle

\noindent\emph{Source and licence.} A sub-Riemannian endpoint map with no regular values. Version arXiv:2609.05991v1 (CC BY 4.0), distributed by arXiv under the Creative Commons Attribution 4.0 International licence, http://creativecommons.org/licenses/by/4.0/, as recorded in the arXiv licence metadata for this exact version and shown on the abstract page https://arxiv.org/abs/2609.05991v1. Full licence notice: \texttt{licenses/arxiv-2609.05991-CC-BY-4.0.txt}.\par\smallskip

\noindent\textit{Editorial scope.} Counted unit: the article's lemma lem:abnormal (source0.tex lines 390--404). The article's complete original proof of it is delivered with it (source0.tex lines 406--522, one \texttt{proof} environment of 117 lines). No mathematics is written, repaired or re-derived: the statement and the proof are the article's own lines, in the article's own notation, and no retained line is substituted or re-flowed.  Retained uncounted as the source-local context the proof needs: lines 373--379.  Nothing was omitted from any retained span, so the package carries no omission record.  No retained line contains a citation, so the wrapper carries no bibliography and no bibliography entry.  No retained line contains a figure environment, an image-inclusion command or drawing code, so no image is delivered and the package declares no asset dependency.  Editorial formatting only, in addition: an AMS \texttt{amsart} preamble that restates the article's own declarations and macros used by the retained lines, namely the environments \texttt{lemma} on separate counters, together with the standard packages those lines need.  The retained environments print in this document as Lemma 1 in order of appearance, whereas the article prints the same items with the numbering of its own full document.  The complete native source of the article as distributed in the arXiv e-print for this exact version is bundled unchanged as \texttt{sources/209481/source0.tex}.
\begin{equation}\label{eq:triangular-system}
 \dot x=u,
 \qquad
 \dot y=v,
 \qquad
 \dot z=-uV(x,y),
\end{equation}

\begin{lemma}[Criterion for singular curves]\label{lem:abnormal}
Let $L>0$ and let $(u,v)\in L^2([0,L],\R^2)$.  Denote by
\begin{equation}
 \gamma=\gamma_{u,v}=(x,y,z):[0,L]\longrightarrow\R^4
\end{equation}
the associated horizontal trajectory starting at the origin.  Suppose that for some 
$\lambda \in\R^2\setminus\{0\}$,
\begin{equation}
 \lambda\cdot F(x(t),y(t))=0
 \qquad\text{for a.e. }t\in[0,L].
\end{equation}
Then $(u,v)$ is a critical point of $\End_0^L$.  Moreover, the
control of the affine reparametrization
$\bar\gamma(t)\coloneq \gamma(Lt)$, $t\in[0,1]$, is critical for $\End_0$.
\end{lemma}

\begin{proof}
By \eqref{eq:triangular-system}, it holds
\begin{equation}
 x(t)=\int_0^t u(s)\dd s,
 \qquad
 y(t)=\int_0^t v(s)\dd s,
\end{equation}
and therefore
\begin{equation}
 \End_0^L(u,v)
 =
 \left(
 x(L),y(L),
 -\int_0^L u(t)V(x(t),y(t))\dd t
 \right).
\end{equation}
Let $(\widehat u,\widehat v)\in L^2([0,L],\R^2)$ be arbitrary, and set
\begin{equation}
 \delta x(t)\coloneq \int_0^t\widehat u(s)\dd s,
 \qquad
 \delta y(t)\coloneq \int_0^t\widehat v(s)\dd s.
\end{equation}
Differentiating the endpoint map in the direction
$(\widehat u,\widehat v)$ gives
\begin{equation}
 D_{(u,v)}\End_0^L(\widehat u,\widehat v)
 =
 \bigl(\delta x(L),\delta y(L),\delta z(L)\bigr),
\end{equation}
where
\begin{equation}
 \delta z(L)
 =
 -\int_0^L
 \left[
 \widehat u V
 +uV_x \delta x
 +uF \delta y
 \right]\dd t.
\end{equation}
Here and below, $V$, $V_x$, and $F$ inside the integrals are evaluated at
$(x(t),y(t))$.

Taking the scalar product with $\lambda$, we obtain
\begin{equation}\label{eq:identityabove}
 \lambda\cdot\delta z(L)
 =
 -\int_0^L
 \left[
 (\lambda\cdot V)\widehat u
 +u(\lambda\cdot V_x)\delta x
 +u(\lambda\cdot F)\delta y
 \right]\dd t.
\end{equation}
Along the reference trajectory,
\begin{equation}
 \frac{\dd}{\dd t}\bigl(\lambda\cdot V(x(t),y(t))\bigr)
 =
 u(\lambda\cdot V_x)+v(\lambda\cdot F)
 \qquad\text{for a.e. }t,
\end{equation}
and
\begin{equation}
 \frac{\dd}{\dd t}\delta x(t)=\widehat u(t)
 \qquad\text{for a.e. }t,
 \qquad
 \delta x(0)=0.
\end{equation}
Hence integration by parts gives
\begin{equation}
 \int_0^L(\lambda\cdot V)\widehat u\dd t
 =
 (\lambda\cdot V(x(L),y(L)))\delta x(L)
 -
 \int_0^L
 \left[
 u(\lambda\cdot V_x)+v(\lambda\cdot F)
 \right]\delta x\dd t.
\end{equation}
Substituting this identity in \eqref{eq:identityabove} yields
\begin{equation}
 \lambda\cdot\delta z(L)
 =
 -(\lambda\cdot V(x(L),y(L)))\delta x(L)
 +
 \int_0^L
 (\lambda\cdot F)
 \bigl(v\delta x-u\delta y\bigr)\dd t.
\end{equation}
By assumption the integral vanishes. Therefore, defining $\eta_L \in \R^4$ by
\begin{equation}
 \eta_L
 \coloneq 
 \bigl(
 \lambda\cdot V(x(L),y(L)), 0, \lambda
 \bigr),
\end{equation}
it holds 
\begin{equation}
 \operatorname{Im}D_{(u,v)}\End_0^L\subset\ker\eta_L.
\end{equation}
Since $\lambda\neq0$, the covector $\eta_L$ is nonzero.  Hence $D_{(u,v)}\End_0^L$ is not surjective.

Finally, define the linear isomorphism
\begin{equation}
 R_L:L^2([0,L],\R^2)\longrightarrow L^2([0,1],\R^2),
 \qquad
 R_L(u,v)(t)\coloneq L\,(u(Lt),v(Lt)).
\end{equation}
The trajectory associated with $R_L(u,v)$ is $t\mapsto\gamma(Lt)$, and
\begin{equation}
 \End_0\circ R_L=\End_0^L.
\end{equation}
Since $R_L$ is an isomorphism, the two endpoint differentials have the
same image.  The reparametrized control is therefore a critical point for
$\End_0$.
\end{proof}

\end{document}
